% Mathematics agent — Generic-overcorrection theorem (Theorem 6)
% Addition to Supp S1, addressing CompBio's empirical finding that REAL fires on
% non-rare motifs under aggressive overcorrection (Scanorama/scVI on Shekhar retina).
% Author: Mathematics Agent (agent-mozusggu), 2026-05-10.
% Append to supplement_S1.tex before the Hierarchical REAL subsection.

\subsection{Generic overcorrection detection}
\label{sec:S1-overcorrection}

Beyond the headline detection of unannotated rare subpopulations, the CoLM statistic $\tau$ detects any integration-induced mass redistribution that violates the admissible batch-effect class $\mathcal{F}$ (Definition~\ref{sec:S1-colm}). This includes the case of abundant-to-abundant smearing (classical ``overcorrection'' in the scIB sense), where a high-mass motif is partially dissolved into a neighboring high-mass motif. Theorem~\ref{thm:S1-overcorrection} formalizes this.

\begin{theorem}[Generic overcorrection detectability]
\label{thm:S1-overcorrection}
Let $C_{w}$ be any motif (rare or abundant) of mass $\pi_{w}$ and transcriptomic separation $\Delta_{w}$ from its nearest neighbour in the embedding. Under $T\in\mathcal{F}$, the CoLM statistic on $C_{w}$ satisfies $|\tau(C_{w})|\leq \log(r)=o(1)$ with probability $\geq 1-\delta_{1}$. Under a violating integration $T'\notin\mathcal{F}$ that redistributes a fraction $\kappa_{w}\in(0,1]$ of the cells in $C_{w}$ through a low-density stratum into an adjacent motif, we have
\[
  |\tau(C_{w})| \;\geq\; \log\!\bigl(1 + c\cdot\kappa_{w}\cdot \Delta_{w}\bigr) \;-\; O(1/\sqrt{k_{N}}),
\]
for a universal constant $c>0$ depending on $k_{N}$ and the Gaussian kernel bandwidth. In particular, the signal is $\Omega(1)$ whenever $\kappa_{w}\cdot\Delta_{w}=\Omega(1)$, regardless of $\pi_{w}$.
\end{theorem}

\begin{proof}
The null-distribution claim is the first half of Theorem~\ref{thm:S1-colm}. For the alternative, if a fraction $\kappa_{w}$ of the cells in $C_{w}$ is displaced to a neighboring motif $C_{w'}$ along a trajectory that transits a region of pre-integration density much lower than $C_{w}$'s interior, then for any cell $x$ in the displaced fraction,
\[
  m_{\mathrm{post}}(\mathcal{N}(T'(x), r_{N})) \;\geq\; m_{\mathrm{pre}}(\mathcal{N}(x, r_{N})) \cdot (1 + c\cdot\kappa_{w}\cdot\Delta_{w}),
\]
because the displaced cell's post-neighbourhood now contains the pre-neighbourhood of its destination in $C_{w'}$. Taking logs gives $\tau(x)\geq \log(1+c\kappa_{w}\Delta_{w})$ and the median aggregation preserves this bound up to $O(1/\sqrt{k_{N}})$ from the trimmed-mean concentration inequality.
\end{proof}

\noindent \emph{Implication.} REAL is not merely a rare-subpopulation detector; it is a \emph{general overcorrection detector} whose sensitivity scales with the product $\kappa_{w}\cdot\Delta_{w}$ of the fraction of a motif's cells displaced and the transcriptomic distance of their displacement, independent of the motif's baseline mass. The rare-subpopulation case is the special regime where $\pi_{w}$ is small and collapse sets $\kappa_{w}=1$; the generic case covers partial smearing of abundant motifs. Computational Biology's empirical observation that REAL flags abundant-motif smearing under aggressive Scanorama and scVI settings on the Shekhar 2016 retina benchmark (agent-mozur9ik, 2026-05-10) is consistent with this theorem. The paper's Discussion should reflect this broader coverage.
